\documentclass[10pt,a4paper]{amsart}
\usepackage[margin=19mm]{geometry}
\usepackage{amsmath,amssymb,mathtools}
\usepackage[scr=rsfs,cal=euler]{mathalfa}
\usepackage{xfrac}
\usepackage[unicode,hidelinks,bookmarks=false]{hyperref}
\newcommand{\ie}{i.e.\ }
\newcommand{\cf}{cf.\ }
\newcommand{\resp}{respectively\ }
\newcommand{\Subsection}{\subsection}
\providecommand{\nobreakdash}{\nobreak\hbox{-}}
\setlength{\emergencystretch}{2em}
\multlinegap0pt
\usepackage{bm}
\usepackage{bbm}
\usepackage{dsfont}
\usepackage{xspace}
\usepackage{cancel}
\usepackage{mathtools}
\usepackage{amsthm}
\makeatletter
\@ifundefined{theorem}{\newtheorem{theorem}{Theorem}}{}
\makeatother
\title[Explicit invariant-preserving integration of differential eq]{Explicit invariant-preserving integration of differential equations using homogeneous projection}
\author{Benjamin Kwanen Tapley}
\date{Source: arXiv:2511.02131v1 (CC BY 4.0)}
\begin{document}
\maketitle

\noindent\emph{Source and licence.} Explicit invariant-preserving integration of differential equations using homogeneous projection. Version arXiv:2511.02131v1 (CC BY 4.0), distributed under the Creative Commons Attribution 4.0 International licence, as recorded in the arXiv licence metadata for this exact version and shown on the abstract page https://arxiv.org/abs/2511.02131v1. Full rights notice: licenses/arxiv-2511.02131-CC-BY-4.0.txt.\par\smallskip

\noindent\textit{Editorial scope.} Counted unit: the Theorem whose source label is \texttt{thm:pseudo-energy} in the article (source0.tex lines 557--566, source label \texttt{thm:pseudo-energy}), delivered together with the article's own complete original proof of it (source0.tex lines 568--611, one \texttt{proof} environment of 44 lines). No mathematics is written, repaired or re-derived: statement and proof are the article's own text line for line. Required context retained uncounted: none. Every definition, display and label of those lines is retained unchanged. Omitted from this package and not redistributed: 1--556, 567--567, 612--1046. Nothing inside a retained excerpt is omitted or altered: every retained body is delivered exactly as the article's lines are written, so no omission receipt of any kind accompanies this package. Citations: the retained excerpts carry no citation command at all, so no bibliography entry of the article is transcribed and no reference is redistributed. Reference adaptation: none was needed and none was made; every reference inside the delivered unit resolves inside it, so no unresolved reference or citation is produced. Printed slips kept literal: no printed word, symbol, hypothesis or number of the counted statement or of its proof was altered. Layout and class: the article's own class is \texttt{elsarticle} with packages including amssymb, amsmath, amsthm, mathtools, bm, subcaption, xcolor, graphicx; the wrapper is the lane's pinned \texttt{amsart} editorial wrapper at 10pt on A4, which restates the theorem-like environments the retained text needs and adds the article's own macro definitions that the retained text needs (none), with extra wrapper packages (bm, bbm, dsfont, xspace, cancel, mathtools). The delivered numbering is the unit's own. Assets: no retained span contains a figure environment, a graphics-inclusion command, a PDF-page inclusion, a table or a TikZ picture; no image file is copied into this package, the package declares no asset dependency and no image file is redistributed with it. Licence: version arXiv:2511.02131v1 of this article is distributed under the Creative Commons Attribution 4.0 International licence, as recorded in the arXiv licence metadata for this exact version and shown on the abstract page https://arxiv.org/abs/2511.02131v1. A full rights notice is delivered at licenses/arxiv-2511.02131-CC-BY-4.0.txt. Characters: every retained excerpt is pure ASCII, so the delivered engine, XeLaTeX with its default Latin Modern text font, reports no missing character.
\begin{theorem}[Pseudo-invariant-preservation]\label{thm:pseudo-energy}
  Let $\Phi_h$ be an order-$p$ one-step method for $\dot x=f(x)$ with exact flow $\varphi_h$ with $m$ invariants $f^T\nabla H_i=0$, such that $H_i(x)H_i(\Phi_h(x))>0$ and each $H_i$ are Lipschitz continuous. Furthermore, let $\Psi_h$ be an order-$q$ one-step method for $\dot z=g(z)$, as defined. Then the method
  $$ \widehat{\Phi}_h(x)\;:=\;\Psi_{1}\circ\Phi_h(x), $$
  where $k_i \;=\; \log\!\left(\frac{H_i(x)}{H_i(\Phi_h(x))}\right)$ for $i=1,...,m$ satisfies the following:
\begin{enumerate}
  \item $|H_i(\widehat{\Phi}_h(x)) - H_i(x)| = \mathcal{O}\big(h^{(p+1)(q+1)}\big)$ for each $i$ (pseudo-invariant-preserving), and 
  \item  $|\widehat{\Phi}_h(x) - \varphi_h(x)\| = \mathcal{O}(h^{p+1})$ (same order as $\Phi_h$).
\end{enumerate}

\end{theorem}

\begin{proof} 

  % \begin{proof}
Define for each $i$ the energy error $\delta_{i,h}:=H_i(\Phi_h(x))-H_i(x)$, which satisfies
$$ |\delta_{i,h}|\le L_{H_i}\,\|\Phi_h(x)-\varphi_h(x)\|=\mathcal O(h^{p+1}).$$
Letting $g$ be defined as above gives $g(z)^\top\nabla H_i(z)=k_i H_i(z)$, by construction, hence along $\psi_t$ we have $H_i(\psi_t(z))=e^{k_i t}H_i(z)$. Evaluating at $z=\Phi_h(x)$, $t=1$ and $k_i \;=\; \log\!\left(\frac{H_i(x)}{H_i(\Phi_h(x))}\right)$ gives 
\begin{equation}\label{eq:pseudo energy exact}
  H_i\big(\psi_1(\Phi_h(x))\big)=e^{k_i}H_i(\Phi_h(x))=H_i(x).
\end{equation}
The order of $\Psi_h$ implies that
$$
\|\Psi_1(z)-\psi_1(z)\|\;=\;\mathcal O\!\big(\|g\|_{C^q}^{\,q+1}\big).
$$
Since $g_i$ is bounded and depends linearly on $k_i$, we have $\|g\|_{C^q}=\mathcal O(|k|)$, hence
$$
\|\Psi_1(z)-\psi_1(z)\|=\mathcal O(|k|^{\,q+1}).
$$
Evaluating this at $z=\Phi_h(x)$ and applying the mean value inequality for $H_i$ gives
$$
\big|H_i(\Psi_1(\Phi_h(x)))-H_i(\psi_1(\Phi_h(x)))\big|
\le L_{H_i}\,\|\Psi_1(\Phi_h(x))-\psi_1(\Phi_h(x))\|
=\mathcal O(|k|^{\,q+1}).
$$
Combining with \eqref{eq:pseudo energy exact} and $|k|=\mathcal O(h^{p+1})$ yields, for each $i$,
$$
\big|H_i(\Psi_1(\Phi_h(x)))-H_i(x)\big|
=\mathcal O\!\big(h^{(p+1)(q+1)}\big),
$$
which proves item 1. For the state error, use the triangle inequality
$$
\|\Psi_1(\Phi_h(x))-\varphi_h(x)\|
\le \|\Phi_h(x)-\varphi_h(x)\|+\|\psi_1(\Phi_h(x))-\Phi_h(x)\|+\|\Psi_1(\Phi_h(x))-\psi_1(\Phi_h(x))\|.
$$
The first term is $\mathcal O(h^{p+1})$ by the order of $\Phi_h$. For the second term,
$$
\|\psi_1(\Phi_h(x))-\Phi_h(x)\|
\le \int_0^1 \|g(\psi_t(\Phi_h(x)))\|\,dt \le C\,|k|=\mathcal O(h^{p+1}),
$$
for some constant $C$. The final term is $\mathcal O(|k|^{\,q+1})=\mathcal O(h^{(p+1)(q+1)})$. Therefore
$$
\|\Psi_1(\Phi_h(x))-\varphi_h(x)\|=\mathcal O(h^{p+1}),
$$
which proves item 2.
\end{proof}

\end{document}
